\subsection{通常型の全共鳴排除と次の無上界問題}
\label{notes:g51522}

\paragraph{第22版で確定した追加結果．}
公開原稿の命題 \texttt{g51522:no-resonance} により，$G_{5,15}$ 通常型十四セルでは，
共有高 ray の $D$ 相殺式の一次係数が零になる共鳴分岐は全て排除された．
正側局所模型と共有 node 分配を固定すれば，負側の残る scale 候補は高々一つである．
これは全 $Q$ に対する定理であり，有界探索の不在を一般化したものではない．
固定した正側の低 row・column の復元は従来の genuine な二次式のままであり，
その全 roots を調べる必要は残る．

\paragraph{再利用する証明の核．}
負側低 row を $(-j,j+m-1,-m)$，正側と負側の共有高 rows を
$(-M,-N,U)$ と $(X,-W,-U)$ とする．
$Q=mU+1$，$M+N=Q+U$，$W=X+Q-U$ であり，共鳴は
$(m-1)X=j(Q-U)$ と同値である．$m=1$ は直ちに不可能，$m=4,j=1$ は
$3X=3U+1$ で不可能である．
$m=2$ では，$X$ に対応する node $h$ から
\[
 M=(j-2+2h)U+j-1+h,\quad
 N=(5-j-2h)U+2-j-h
\]
となる．$j\ge3$ と $N>0$ は $h=0,j=3,4$ を強制する．
共鳴時の負側共有係数は，$J=j(j+1)$ として
\[
 D_{\rm sh}^-=-2d_--2-2/J+
 \{J(d_-+2)+2\}/\{J(U+1)\}.
\]
整数性から $(j,d_-,U)=(3,0,12),(4,0,20),(4,5,70)$ に限られ，
正側の共有 pair の整除性が全てを排除する．$M_2$ にある二つの負成分の交換も含む．

\paragraph{退化する $T=1$ の扱い．}
五つの固定低データ問題を低 row $(-1,-m,m)$，低 $D=d$，
低 column $(-S,-T,U)$ で表す．
公開原稿の命題 \texttt{g51522:T1} は全五問題で $T=1$ を排除する．
$m=1$ は第二分母が二となるため，$m=2,4$ は元の row/pair equations からの
\[
 K=\frac{\{(2m+1)S+m\}U+S+m}{(mS-1)U+S-1},\qquad U=S-d
\]
による．$K>2$ と $K\ge3$ の整数性は $m=2$ なら $S=2,4$，
$m=4$ なら $S=2$ を強制し，残りを厳密な整除条件で排除する．
この境界を無視して二次式の先頭係数で割ってはならない．

\paragraph{今後の有用な注意．}
canonical Schur の向き $k\ge0$ と，低 row の順序 $(-1,-m,m)$ は独立に固定できない．
二つの負の標識を交換すると $k'=-k-1$ となるので，
Schur の符号定理を用いる際は $(s,t)=(1,m),(m,1)$ の両向きを覆う．
本版の共鳴と $T=1$ の証明はこの向きの選択を使っていない．
また，低 $D$ の負性だけから模型を排除してはならない．
例えば低 row $(-1,-1,1)$，低 column $(-4,-58,292)$ には
matching entries $H=7,K=42$ を持つ正の奇分母整数模型が存在し，
低 $D=-231$ となる．これは対象 $d=-9$ の反例ではなく，
一般の低 $D\ge0$ という推測を避けるための小さな反例である．

\paragraph{次の研究課題．}
非共鳴の五つの固定 $(\text{low row},d)$ 問題は全 $Q$ に対して未分類である．
新しい上界を持たない $Q$ の scan を増やすより，
共有 $D$ の一次解を負側の残る pair equation に代入し，
正側の二次式との消去から scale の上界または符号障害を得ることを優先する．
特に $T\ge2$ に限定できたので，次の元の恒等式が利用できる．
\[
 A=(S+1)(mT+1)-1,\quad B=(m-1)T+1,\quad
 C=(S+1)(mU+1)-1,
\]
\[
 E=(m+1)U+1,\quad F=(mT+1)(mU+1)-1,
\]
\[
 (F-EB)H^2+\{CB+F-E(A+B)\}H+A(C-E)=0,
 \qquad F-EB=(T-1)(U+1)>0.
\]
ここで $U=S+T-d-1$，$H>0$，$K=A/H>0$ であり，
$H\mid A$，$H+1\mid C$，$K+B\mid F$ と元の row sum を課す．
この表示は一つの消去変数を減らすが，現時点で全 $Q$ の有限性を証明するものではない．
負例外 $Q=3,5$ も未解決として残す．

\paragraph{検算と再現．}
\nolinkurl{verify_g5_15_resonance_boundary_v22.py} は旧版ファイルを import せず，
係数恒等式，十四セルの網羅性，証明で得た定数 $26,42,142$ の全正約数，
五つの順序付き pair 検査，五つの $T=1$ 境界を標準ライブラリで検算する．
結果は \nolinkurl{g5_15_resonance_boundary_verification_v22.json} に記録される．
この検算は証明支援系での形式化でも外部査読でもない．
